\section{Second-order plasmonic response}
\label{app:perturbation}

In units of the contact width $W$, the plasma speed $s_{0}$ and the gated
density $n_{0}$, Eqs.~\eqref{eq:cont} and \eqref{eq:mom} contain the
friction $\Gamma=W/s_{0}\taumr$, the plasma Reynolds number
$\Rey_{s}=s_{0}W/\nu$, the collision time $s_{0}\tauee/W$ and the
quantum-capacitance fraction $q$, with the closure
$h(n)=2(1-q)(\sqrt{n}-1)+q\ln n$ and $h''(1)=-(1+q)/2\equiv\kappa$. Writing
$n=1+\mathrm{Re}(n_{1}e^{i\Omega t})+n_{20}$ and
$\bm{u}=\mathrm{Re}(\bm{u}_{1}e^{i\Omega t})+\bm{u}_{20}$, the first order is
\begin{align}
i\Omega n_{1}+\nabla\cdot\bm{u}_{1}&=0,\\
(i\Omega+\Gamma)\bm{u}_{1}+\nabla n_{1}-\nu(\Omega)\nabla^{2}\bm{u}_{1}&=0,
\end{align}
with $\nu(\Omega)=\Rey_{s}^{-1}/(1+i\Omega s_{0}\tauee/W)$, and the time
average of the second order is
\begin{align}
\nabla\cdot\bm{u}_{20}&=-\tfrac12\mathrm{Re}\,\nabla\cdot(n_{1}\bm{u}_{1}^{*}),\\
\Gamma\bm{u}_{20}+\nabla P-\Rey_{s}^{-1}\nabla^{2}\bm{u}_{20}
&=-\tfrac12\mathrm{Re}\,(\bm{u}_{1}^{*}\cdot\nabla)\bm{u}_{1},
\end{align}
where $P=n_{20}+(\kappa/4)|n_{1}|^{2}$. The dc contact voltage is
$P+c_{V}|n_{1}|^{2}$ with $c_{V}=1/8$ for the Dirac closure and zero for
shallow water, and the dc particle current through a contact is the flux of
$\bm{u}_{20}+\mathrm{Re}(n_{1}\bm{u}_{1}^{*})/2$. Each contact is an
equipotential with one unknown voltage per order, closed by Kirchhoff's law
for its termination: $R I_{c}+V_{c}=V_{a,c}$ for a source of resistance $R$,
$I_{c}=0$ for an open contact. The same structure with a matched load at dc
gives the rectified power. Voltages are in units of
$V_{u}=m^{*}s_{0}^{2}/e$, resistances in units of
$m^{*}s_{0}/e^{2}n_{0}W$ ($65\ \Omega$ for $W=4\um$), and frequency in units
of $s_{0}/2\pi W$ ($0.142\THz$).

A contact resistance $R_{c}$ enters as a change of the ports only: the
source becomes a Thevenin source of resistance $\Rant+R_{c}$ and the drain
one of zero amplitude and resistance $R_{c}$. In second order the dc
problem is linear, so the short-circuit current and open-circuit voltage
satisfy $I_{\mathrm{sc}}(\Rdc+2R_{c})=V_{\mathrm{oc}}$ exactly, which is
one of the ladder tests. At the signal frequency the antenna does not see
the device impedance plus $2R_{c}$, because the source and drain currents
differ by the gate displacement current.

\section{Saturation estimate on the density and gate grid}
\label{app:knee}

The peak efficiencies of Fig.~\ref{fig:operating}(b) are
$\eta_{\mathrm{peak}}=2r^{2}(\Rant/\Rdc)\,(0.7)^{2}\,\mathcal{S}$, with $r$
and $\Rdc$ from the small-signal solver at each grid point and
$\mathcal{S}$ the ratio of the time-domain efficiency to the square law at
$V_{a}=0.7\,V_{u}$ for the reference device ($0.86$, $0.89$, $0.60$ at
$100$~K and $1.14$, $1.10$, $0.81$ at $300$~K for $0.1$, $0.3$, $1\THz$).
The assumption that saturation occurs at a fixed $V_{a}/V_{u}$ was tested
by a time-domain ladder at $(5\times10^{11}\percm,10\nm)$: $\mathcal{S}$
agrees to $0.2\%$ at $300$~K and is $0.76$ instead of $0.89$ at $100$~K, so
the $100$~K peak values at low density and thin gates are upper estimates
by $10$ to $20\%$.

\section{Steady gated flow in one dimension}
\label{app:gated}

In a gradual channel of width $W(x)$ carrying the current $I$ the
continuity equation gives $n u W=I$. Without the convective and viscous
terms the steady form of Eq.~\eqref{eq:mom} is $\Gamma u=-h'(n)\,dn/dx$ in
the units of Appendix~\ref{app:perturbation}, so that
\begin{equation}
n\,h'(n)\,\frac{dn}{dx}=-\frac{\Gamma I}{W(x)} .
\end{equation}
With $G'(n)=n h'(n)$, i.e. $G=n^{2}/2$ for shallow water and
$G=\tfrac23(1-q)n^{3/2}+qn$ for the Dirac closure, integration from the
driven contact, where $n=n_{s}$, to the contact tied to the gate, where
$n=1$, gives Eq.~\eqref{eq:1d}. The voltage of the driven contact is
$\Phi(n_{s})$ with $\Phi=(1-q)(n-1)+2q(\sqrt{n}-1)$ in units of $V_{u}$.
The width enters only through $\int dx/W$, which is unchanged when the
channel is turned end for end, so exchanging the contacts cannot produce
an asymmetry. Reversing the current changes the sign of the right-hand
side; expanding to second order with $G''(1)=(1-q)/2$ and
$\Phi''(1)=-q/2$,
\begin{equation}
V=\rho I-\tfrac14(\rho I)^{2},\qquad \rho=\Gamma\int_{0}^{L}\frac{dx}{W},
\end{equation}
and the ratio of the two voltage magnitudes is $1+\rho|I|/2=1+V/2V_{u}$,
which is Eq.~\eqref{eq:dpol} with $U_{0}=(1-q)V_{u}$ and $q=0.039$; for
shallow water the coefficient of $(\rho I)^{2}$ is $\tfrac12$ and the ratio
is $1+V/U_{0}$. In a uniform channel the convective term can be kept: the
equation becomes $(nh'-j^{2}/n^{2})\,dn/dx=-\Gamma j$ with $j=I/W$, with the
integral $G(n)+j^{2}/n$, and the two-dimensional solver reproduces this
closed form to $1.1\times10^{-4}$ for both closures and both signs up to
$V=0.3\,U_{0}$. In a channel of varying width the convective term brings in
$dW/dx$, and with it the inertial shape term of Sec.~\ref{sec:dc}, which is
second order in the current.

